\documentclass{article}

\newcommand{\svninfo}{Retistruct version: 0.5.11 , Date: 2016-06-29}
\pagestyle{myheadings}
\markboth{\svninfo}{\svninfo}

\usepackage[a4paper,left=1in,right=1in,top=1in,bottom=1in,head=1in]{geometry}

\usepackage[british]{babel}
\usepackage{tabularx}
\usepackage{graphicx}
\usepackage{hyperref}

\title{\textsc{Retistruct} manual}
\author{David C. Sterratt}

\begin{document}

\maketitle
\thispagestyle{myheadings}

\section{Installation}
\label{manual:sec:installation}

\subsection{Install the necessary system packages}

\subsubsection{Fedora Linux }

As root,

\begin{verbatim}
yum install R-core
yum install R-devel
yum install gtk2
yum install gtk2-devel
yum install mesa-libGL-devel
yum install mesa-libGLU-devel
\end{verbatim}

\subsubsection{Ubuntu Linux}
At the command line, enter:

\begin{verbatim}
sudo apt-get install r-base r-cran-rgl libgtk2.0-dev
\end{verbatim}
% Need to install RGtk2 separately??

\subsubsection{Mac OS X}
\label{retistruct-manual:sec:mac}

First install R for Mac OS X, available from
\url{http://www.r-project.org/}. Under Mac OS X Leopard and Lion, this
should be sufficient. 

\paragraph{Under Mac OS X 10.8 (Snow Leopard),} it may be necessary to
install the following external packages:
\begin{itemize}
\item Xcode (free from App Store)
\item Xquartz (\href{http://xquartz.macosforge.org/}{http://xquartz.macosforge.org/})
\item GTK (\href{http://r.research.att.com/libs/GTK_2.18.5-X11.pkg}{http://r.research.att.com/libs/GTK\_2.18.5-X11.pkg})
\end{itemize}

\paragraph{Under Mac OS X 10.9 (Mavericks) or Mac OS X 10.10
  (Yosemite),}
follow Chris von Csefalvay's instructions:
\href{http://chrisvoncsefalvay.com/2015/02/08/installing-gtk-on-mavericks.html}{http://chrisvoncsefalvay.com/2015/02/08/installing-gtk-on-mavericks.html}.

\subsubsection{MS/Windows}

First install R for MS/Windows, available from \url{http://www.r-project.org/}.

\subsection{Install the core \textsc{Retistruct} package}

\begin{enumerate}
\item Start \textsc{R}
 \item  Type:
\begin{verbatim}
install.packages("retistruct")
\end{verbatim}
   The first time this runs, it should create a personal directory for
   R packages, and it will take a few minutes to install some required
   packages.
 \item Test that it is working as follows (from within R)
\begin{verbatim}
   library(retistruct)
   demo("retistruct.method")
\end{verbatim}
   If all is well, you should be prompted to press \textbf{Return} a
   few times and Figure 1 from the paper \cite{SterrattEtal2012} will
   appear (drawing the 3D renditions takes some time). There will be
   reports of warnings; they are to do with printing the figures to
   files. You can view them by typing \texttt{warnings()}.
 \item A further test produces Figure~2:
\begin{verbatim}
   demo("low.high")
\end{verbatim}
 \item A further (quite long-running) test produces Figure~6:
\begin{verbatim}
   demo("figure6")
\end{verbatim}
\end{enumerate}

\subsection{Install the \textsc{Retistruct} GUI}
\label{retistruct-manual:sec:inst-retistr-gui}

If the previous step works:
\begin{enumerate}
\item Type:
\begin{verbatim}
retistruct()
\end{verbatim}
\item Some more packages will be downloaded -- this may take a
  minute or so.
\item If you are using MS/Windows, there may be an error message
  saying that libatk-1.0-0.dll is missing from your computer. Click on
  \textbf{OK}. Another message will appear, asking if you want to
  install GTK+. Agree to this.  After this has finished, quit R
  (without saving the workspace image) and restart.
\item If all works, an interface window should appear. A number of
  sets of demonstration data are available from the \textbf{Demo} menu item.
\end{enumerate}

% \section{A tutorial}
% \label{retistruct-user-guide:sec:tutorial}

% This is a tutorial using some sample data provided on the website. 

% \begin{enumerate}
% \item Download the data from
%   \href{http://retistruct.r-forge.r-project.org/Figure_6-data.zip}{http://retistruct.r-forge.r-project.org/Figure\_6-data.zip}. This
%   is the data that underlies Figure~6 of \cite{SterrattEtal2012}
% \item Unpack the zip file. A directory \texttt{Figure\_6-data} will be created.
% \item Start retistruct by first starting \textsc{R}. At the \textsc{R}
%   prompt type:
% \begin{verbatim}
% > library(retistructgui)
% > retistruct()
% \end{verbatim}
%   A window should appear.
% \item Click on the \textbf{Open} icon and navigate to the
%   \texttt{left-contra} directory in the \texttt{Figure\_6-data} directory.
% \end{enumerate}

\section{Running \textsc{Retistruct}}
\label{manual:sec:running}

To start the program, start \textsc{R}. At the \textsc{R} prompt type:

\begin{verbatim}
library(retistruct)
retistruct()
\end{verbatim}
A window should appear.

\subsection{Opening the files for a retina}
\label{manual:sec:opening-files-retina}

There are a number of types of information associated with a
flat-mount retina that \textsc{Retistruct} can process:
\begin{itemize}
\item The coordinates of the outline of the flattened retina
\item An image of the flat-mount retina (optional)
\item The coordinates of labelled data points within the flat-mount
  retina (optional)
\end{itemize}
To import this information into \textsc{Retistruct}, for each retina,
a directory should be created containing files with the above
information, as will be described below.  At present there are three
formats of directory that \textsc{Retistruct} can read.  Most users
will probably find the the \textsc{ImageJ} ROI format most convenient.

\subsubsection{\textsc{ImageJ} ROI format}
\label{retistruct-manual:sec:ijroi-format}

This format allows you to load images of retinae whose outlines have
been marked up in \textsc{ImageJ}. To create the image and outline
files for this format:
\begin{enumerate}
\item Create a directory to save the files created to.
\item Open up \textsc{ImageJ} (or FIJI).
\item Use \textbf{File$\rightarrow$Open} to open the image.
\item Use \textbf{Image$\rightarrow$Scale} to down sample so that the
  resolution is less than 1000x1000. (This is not crucial, but will
  speed up things later.)
\item Save the down-sampled version of the image to the file name
  \texttt{image.png} in the directory created in Step 1.
\item Use the Polygon Tool
  (\includegraphics[height=\baselineskip]{poly}) to mark the edge of
  the retina. According to the \textsc{ImageJ}
  manual\footnote{http://rsbweb.nih.gov/ij/docs/tools.html}: ``To
  create the selection, click repeatedly with the mouse to create line
  segments. When finished, click in the small box at the starting
  point (or double-click), and ImageJ automatically draws the last
  segment.

  The points that define a polygon selection can be moved or deleted,
  and new points can be added. To delete a point, click on it with the
  alt key down. To add a point, click on an existing point with the
  shift key down.''
\item Open the ROI manager by selecting
\textbf{Analyze$\rightarrow$Tools$\rightarrow$ROI Manager}.
\item Click on the \textbf{Add [t]} button.
\item Click on the \textbf{More$\rightarrow$Save} button. In the \textbf{Save
selection\dots} box that appears, enter the file name
\texttt{outline.roi} and make sure this file is saved to the same
directory as the \texttt{image.png} file, i.e.\ the directory created
in Step 1.
\item Optionally mark up the optic disc in the same way as the
  outline. Save this ROI to a file called \texttt{od.roi} in the same
  directory as \texttt{image.png}.
\item Open \textsc{Retistruct}.
\item Click on the \textbf{Open File} icon and select the directory
containing \texttt{image.png} and \texttt{outline.roi}. The retinal
image should now appear, with the outline shown. By default the
outline is in black. If this isn't visible against the image, press
the \textbf{Properties} button in the interface (or select \textbf{Edit$\rightarrow$Properties}), and change the
\textbf{Outline colour}.
\end{enumerate}

\paragraph{Specifying the scale} Optionally, the scale of the image
can be specified by providing a file \texttt{scale.csv} in the same
directory as the image and ROI files. This file should contain two
lines, the first with the heading \texttt{Scale} and the second with
the length of the side of a pixel in micrometres. For example:
\begin{verbatim}
"Scale"
1.5
\end{verbatim}

\paragraph{Reading in data points}
\label{retistruct-user-guide:sec:read-data-points}

In general, the coordinates of data points are read in from a csv file
called \texttt{datapoints.csv}. The two cells of the first line of the
file contain the name of the group of data points and the colour that
these should be displayed in \textsc{Retistruct}.  Marking up the
points from the image in \textsc{ImageJ} will ensure that the
coordinates are in the same system as those used for the image and the
outline. To do this:
\begin{enumerate}
\item Open the image used for marking up the outline
  (\texttt{image.png}) in \textsc{ImageJ}.
\item Threshold the image
  (\textbf{Image$\rightarrow$Adjust$\rightarrow$Threshold\dots} so
  that the labelled points are visible.
\item Select \textbf{Analyze$\rightarrow$Set Measurements\dots} and make
  sure that \textbf{Centroid} is checked
\item Select \textbf{Analyze$\rightarrow$Analyze Particles}. Make sure
  that \textbf{Display results} and \textbf{Clear results} are
  checked. Click on \textbf{OK}.
\item A window entitled \textbf{Results} should appear. This contains
  the X and Y coordinates of the detected points. In this window
  select \textbf{File$\rightarrow$Save as\dots} and save the file as
  \texttt{datapoints.csv}.
\item Open this file with either \textsc{Excel} or \textsc{OpenOffice
    Calc} and remove all the columns apart from the X and Y
  columns. Replace the X with the name of the data set and the Y with
  the colour it should be displayed in \textsc{Retistruct}.
\item Save \texttt{datapoints.csv} in the retina's directory, along
  with \texttt{outline.roi} (and \texttt{image.png} if an image is
  desired).
\end{enumerate}

\subsubsection{CSV format}
\label{retistruct-manual:sec:csv-format}

This is the same as the \textsc{ImageJ ROI} format, except that the
outline is contained in the two columns of a file called
\texttt{outline.csv}.  Each column should have a heading, e.g.:
\begin{verbatim}
"X", "Y"
1,5
10,76,
...
\end{verbatim}
The optic disc outline can be supplied in a similar file called
\texttt{od.csv}.

\subsubsection{IDT format}
\label{retistruct-manual:sec:idt-format}

This is the format used in Ian Thompson's lab (see
Appendix~\ref{manual:sec:reading-data}). These files are contained in
a directory. To open the files corresponding to a retina, click on the
open file icon, and navigate to the directory containing the
\texttt{SYS} and \texttt{MAP} files. On opening this directory, the
retinal outline should appear in the \textsc{Retistruct} window.


\subsection{The retinal display}
\label{manual:sec:retinal-display}

In the ``Show'' section at the bottom left of the screen there are
checkboxes that allow you to show various types of information:
\begin{description}
\item[Markup] Locations of tears and the dorsal or nasal pole
\item[Stitch] Locations of how the algorithm has stitched tears (only
  visible after the reconstruction step)
\item[Grid] Lines of latitude and longitude projected back onto the
  flattened retina (only visible after the reconstruction step)
\item[Datapoints] Locations of data points, such as the locations of
  beads of dye
\item[Landmarks] Landmarks such as the optic disc
\item[Strain] This shows information about how the retina has been
  reconstructed, after the reconstruction step has taken place.
\item[Preserve area] When this is selected, the polar plot displays an
  area-preserving projection
\item[Contours] Display Kernel Density Estimate contours of the data
  points. By default, contours at 5\%, 25\%, 50\%, 75\% and 95\% of the maximum
  density are shown. This can be changed by a command like this at the
  R prompt:
\begin{verbatim}
options(contour.levels=c(5, 25, 50))
\end{verbatim}
  This command would cause the 5\%, 25\% and 50\% contour lines to be
  displayed.
\end{description}

\subsection{Marking up the retina}
\label{manual:sec:marking-up-retina}

\begin{description}
\item[Add tear] To add a tear, click on this button, then click on
  three points in turn which define a tear. The order in which the
  points are added does not matter. Tears contained within a tear can
  be marked up, but tears cannot cross over one another.
\item[Move Point] To move one of the points defining a tear, click on
  this button, then click on the point which you desire to move, then
  click on the point to which it should be moved.
\item[Remove tear] To remove a tear, click on this button, then click
  on the apex of the tear (marked in cyan on the plot)
\item[Mark nasal] To mark the nasal pole, click on this button, then
  click on the point which is the nasal pole.  If the nasal or
  dorsal pole has already been marked, the marker is removed from
  the existing location. The nasal pole should not be in a tear. If
  the nasal tear is placed within a tear, no error is reported at this
  stage, but it will be reported later.
\item[Mark dorsal] As above, except for the dorsal pole.
\item[Mark OD] To mark the optic disc, click on the structure marked
  in orange which you think is the OD. Once clicked on, the structure
  should become blue.
\item[Phi0] This determines the latitude of the rim of the
  reconstructed hemisphere. It depends on the age of the animal.
\end{description}

\subsection{Metadata}
\label{retistruct-manual:sec:metadata}

\begin{description}
\item[Data/Flip DV] Flip the DV axis to compensate for microscope
  orientation. Affects display of retinae.
\item[Eye] Specify whether the eye is Right or Left.  Affects display
  of retinae.
\end{description}

\subsection{Saving the markup and metadata}
\label{manual:sec:saving-markup}

To save the markup, click on the ``Save'' button in the toolbar. This
saves various markup files to the directory containing the
data files. This saved data can be used to reconstruct the retina using
a batch process (Section~\ref{manual:sec:runn-batch-reconstr}).

\subsection{Reconstructing the retina}
\label{manual:sec:reconstr-retina}

To reconstruct the retina, click on the ``Reconstruct retina''
button. This causes a (lengthy) sequence of operations to be performed:
\begin{description}
\item[Stitching] Links between corresponding points on parts of the retinal
  outline  contained in tears are made.
\item[Triangulation] A triangular mesh is placed over the flattened retina
\item[Initial projection to sphere ] The mesh is projected roughly
  onto a sphere
\item[Optimisation] The locations of the mesh points on the hemisphere
  is adjusted so as to minimise a weighted sum of the squared
  differences between the lengths of links in the mesh on the
  hemisphere and on the flattened retina, whilst ensuring that as few
  triangles as possible are flipped.
\end{description}

At the end of the reconstruction process, a polar plot appears next to
the flattened retina. When ``Landmarks'' are shown, the location of
the cuts and tears in the polar coordinates can be seen.

When ``Strain'' is shown, the polar plot is replaced by a scatter plot of the
length of links in the reconstructed object versus the length on the
flattened object. The colours of the points indicate the degree of
expansion or compression from the flattened object to the
reconstructed object.

\subsection{Saving the reconstruction}
\label{manual:sec:saving-reconstr}

To save the markup, click on the ``Save'' button in the toolbar. This
saves various files to the directory containing the data files. When
the files in the directory are opened again, all the markup
information (cut locations, location of nasal point and rim angle) is
loaded, and the reconstruction will also appear, unless there has been
a major upgrade of the software, in which case the retina will need to
be reconstructed using the ``Reconstruct retina'' button.

\subsection{Reading reconstruction data into R}
\label{retistruct-user-guide:sec:reading-date-into}

It can be useful to have access to the data underlying a
reconstruction directly, for example to allow statistical analysis.
To load saved date into R, type the following into R:
\begin{verbatim}
r <- retistruct.read.recdata("/path/to/reconstruction/directory")
\end{verbatim}
The resulting object \texttt{r} contains various fields, which can be
accessed using the R \texttt{\$} operator, e.g.
\begin{verbatim}
r$Dss$green
\end{verbatim}
gives a list of the spherical co-ordinates of green labelled
datapoints. See Table~\ref{tab:matlab-export} for a list of the most
useful quantities. Other quantities (for example the locations of the
mesh points in the flat and spherical retina) are also available. The
documentation for \texttt{ReconstructedOutline},
\texttt{ReconstructedDataset}, \texttt{getDss} and \texttt{getSss} in
the user manual has more information about this.

\begin{table}[p]
  \begin{tabularx}{\linewidth}{p{1.2in}X}
    \hline \texttt{phi0} & The latitude of the rim, expressed in
    radians
    \\
    \texttt{Dss.green}, \texttt{Dss.red} & Locations of green (red)
    labelled cell bodies in spherical coordinates. The first column
    contains the latitude of each point, measured in radians. The
    second column contains the
    longitude of each point, measured in radians. \\
    \texttt{DssMean.green}, \texttt{DssMean.red} & Location of Karcher
    mean of green-labelled cell bodies in spherical coordinates. The
    first column contains the latitude of each point, measured in
    radians. The second column contains the
    longitude of each point, measured in radians. \\
    \texttt{DssHullarea.green}, \texttt{DssHullarea.red} & Area of
    convex hull of green (red) points on sphere. The convex hull is
    essentially a polygon drawn around data points. \\
    \texttt{Tss} & A structure in which each element contains
    spherical coordinates (in the same latitude-longitude format as
    above) of a tear. \\
    \texttt{Sss} & A structure in which each element contains
    spherical coordinates (in the same latitude-longitude format as
    above) of a landmark. \\
    \texttt{KDE} & An object containing information about Kernel
    Density Estimates of the locations of cell bodies.\\
    & \begin{tabular}{p{1.5in}p{3in}} \texttt{green\_flevels} &
      Contour heights, determined by finding heights that exclude a
      certain fraction of the probability. For example, the 95\%
      contour is excludes 95\% of the probability mass,
      and it should enclose about 5\% of the points. \\
      \texttt{green\_labels} & Contour labels. These give the label
      (e.g.\ 5, 25, 50, where these are the percentages above) of each
      contour. Note that there may be more than one contour at the
      same level, so this vector may contain more elements than
      \texttt{flevels}. The first element of \texttt{green\_labels}
      labels the contour whose coordinates are specified in
      \texttt{green\_contours1}, the second element of
      \texttt{green\_labels} relates to \texttt{green\_contours2} and
      so on. \\
      \texttt{green\_tot\_contour\_areas} & The total area in square
      degrees enclosed by each contour. This is a matrix with the
      first column giving the contour label (see above) and the
      next column giving the area. \\
      \texttt{green\_kappa} & The concentration parameter of
      the Fisher density determined by the kernel fitting algorithm. \\
      \texttt{green\_h} & A pseudo-bandwidth parameter, the inverse
      of the square root of \texttt{kappa}. Units of degrees. \\
      \texttt{green\_maxs\_phi} & Lattitude of maximum point of kernel
      estimate. \\
      \texttt{green\_maxs\_lambda} & Longitude of maximum point of kernel
      estimate. \\
      \texttt{green\_g\_xs} \texttt{green\_g\_ys} \texttt{green\_g\_f}
      & Kernel density estimates on standard polar grid. This can be
      plotted in MATLAB using the command
      \texttt{contour(KDE.green\_g\_xs, KDE.green\_g\_ys,
        KDE.green\_g\_f)}. \\
      \texttt{green\_gpa\_xs} \texttt{green\_gpa\_ys}
      \texttt{green\_gpa\_f} & Kernel density estimates on
      area-preserving polar grid. Plotted in MATLAB as above. \\
      \texttt{green\_contours1}, \texttt{green\_contours2\dots} &
      Coordinates of contours. See \texttt{green\_labels} above for
      more explanation. \\
      \texttt{green\_contour\_areas1},
      \texttt{green\_contour\_areas2\dots} & Area contained within each
      individual contour. See \texttt{green\_labels} above for more
      explanation.
    \end{tabular} \\
    \texttt{KR} & An object containing Kernel Regression estimates of the
    density of points, derived from the grouped data points. All
    fields correspond to \texttt{KDE} above.\\
    \hline
  \end{tabularx}
  \caption{Variables exported in the \texttt{r.mat} file.}
  \label{tab:matlab-export}
\end{table}


\subsection{Reading reconstruction data into \textsc{Matlab}}
\label{retistruct-manual:sec:export-reconstr-data}

By default when a reconstruction is saved, a subset of the data is
stored in a file called \texttt{r.mat} in the same directory as the
raw data and the markup. To import this data into \textsc{Matlab}, cd
into that directory, and type:
\begin{verbatim}
clear
load r.mat
\end{verbatim}
This puts a number of variables into the workspace, as shown in
Table~\ref{tab:matlab-export}.

To produce a polar plot of data points, try the following code:
\begin{verbatim}
polar(Dss.green(:,2), Dss.green(:,1)*180/pi+90, '.g')
hold on
polar(Dss.red(:,2),   Dss.red(:,1)*180/pi  +90, '.r')
hold off
\end{verbatim}
The radial axis indicates the latitude in degrees measured from the
retinal pole.

In the \texttt{matlab} subdirectory of the distribution, there are
some scripts to produce polar polar plots, including the locations of
tears and landmarks. To create PDF plots of all the retinae in a
directory, try:
\begin{verbatim}
makefigures('retinae', 'output_directory')
\end{verbatim}

There are also some embryonic scripts to create polar plots:
\texttt{plot\_datapoints\_polararea.m} and \texttt{polararea.m}.




\subsection{Running a batch of reconstructions}
\label{manual:sec:runn-batch-reconstr}

The \textsc{Retistruct} library can be used to reconstruct a batch of
retinae which have been marked up. Suppose that the directory
\texttt{retinae} contains a directory tree in which there are data
directories containing the raw outline, datapoint and image files and
the saved markup files. In order to perform the reconstructions, we
create a new directory \texttt{retinae/reconstructions}, and run the
following sequence of commands in R:
\begin{verbatim}
R
> library(retistruct)
> retistruct.batch(tldir='retinae', outputdir='retinae/reconstructions')
\end{verbatim}
This command will go through the \texttt{retinae} directory, looking
for valid data directories. If it finds one, it sets about trying to
reconstruct the retina. As it reconstructing each retina, it writes to
log file in \texttt{retinae/reconstructions}. Once the reconstruction
is complete, it saves a number of plots in this directory in PDF
format. It also adds a line to a summary log file in
\texttt{retinae/reconstructions} called
\texttt{retistruct-batch.csv}. This file contains a number of columns:
\begin{description}
\item[\texttt{Dataset}] The directory of the data set
\item[\texttt{Return}] The return value from the process
\item[\texttt{Result}] A summary of the result, including if any
  errors were returned
\item[\texttt{E}] The total error  of the optimised reconstruction
\item[\texttt{El}] The error due to purely to the lengths of links in
  the optimised reconstruction
\item[\texttt{nflip}] The number of flipped triangles
\item[\texttt{EOD}] The distance of the Optic Disc from the inferred
  centre of the retina, in degrees. If the OD has not been marked up,
  this is \texttt{NA}.
\end{description}

To export the reconstruction data in a directory hierarchy in which
\texttt{retistruct.batch()} has been run, run the following sequence
of commands in R:
\begin{verbatim}
R
> library(retistruct)
> retistruct.batch.export.matlab(tldir='retinae')
\end{verbatim}


\appendix

\section{The IDT Data format}
\label{manual:sec:reading-data}

\begin{table}
  \begin{tabular}{ll}
    \hline
    \multicolumn{2}{c}{\textbf{FOR EACH BOUNDARY}} \\
    \hline
    \texttt{MAPNUM}   & id number of boundary \\  
    \texttt{MINLAT}   & min latitude      \\
    \texttt{MAXLAT}   & max latitude      \\
    \texttt{MINLON}   & min longitude     \\
    \texttt{MAXLON}   & max longitude     \\
    \texttt{LABLAT}   & latitude of label \\
    \texttt{LABLON}   & longitude of label\\
    \hline
    \multicolumn{2}{c}{\textbf{FOR EACH CELL}} \\
    \hline
    \texttt{XRED}     & $x$-coordinate if cell labelled red but not doubly\\
    \texttt{YRED}     & $y$-coordinate if cell labelled red but not doubly\\
    \texttt{XGREEN}   & $x$-coordinate if cell labelled green but not doubly\\
    \texttt{YGREEN}   & $y$-coordinate if cell labelled green but not doubly\\
    \texttt{XDOUBLE}  & $x$-coordinate if cell labelled doubly\\ 
    \texttt{YDOUBLE}  & $y$-coordinate if cell labelled doubly\\
    \texttt{XGRID}    & sample box cell is in \\
    \texttt{YGRID}    & sample box cell is in \\
    \texttt{PERIM}    & perimeter of cell \\
    \texttt{AREA}     & area of cell \\
    \hline
    \multicolumn{2}{c}{\textbf{ONE PER GRID BOX}} \\
    \hline
    \texttt{GRIDX}    & grid location of centre of sample box \\
    \texttt{GRIDY}    & grid location of centre of sample box \\
    \texttt{XGRIDCOO} & $x$-coordinate of centre of sample box \\
    \texttt{YGRIDCOO} & $y$-coordinate of centre of sample box \\
    \texttt{BOXSIZEX} & size (half width) of sample box in $x$-direction \\
    \texttt{BOXSIZEY} & size (half width) of sample box in $y$-direction \\
    \texttt{COMPLETE} & whether counting of sample has been completed\\
    \texttt{TOTALCEL} & total number of cells in this box\\
    \texttt{TOTALRED} & total number of red-only cells in this box\\
    \texttt{TOTALGRE} & total number of green-only cells in this box\\
    \texttt{TOTALDOU} & total number of double cells in this box\\
    \texttt{MEANPERI} & average perimeter of all cells \\
    \texttt{MEANAREA} & average area of all cells \\
  \end{tabular}
  \caption{Column headings of the \texttt{SYS.SYS} file.}
  \label{tab:data-format}
\end{table}

The data for each retina is stored in a separate directory. Within
each directory there are two files:
\begin{description}
\item[\texttt{SYS.SYS}] A table in SYSTAT format containing the
  coordinates of the red, green and doubly labelled cell bodies, and
  counts of labelled cell bodies within each grid box. The column
  headings shown in Table~\ref{tab:data-format}.  Each row of the
  table contains information only on a subset of the data, e.g.\ the
  coordinates of a red-labelled cell.
\item[\texttt{ALU.MAP}] A text file containing the coordinates of the
  map outline. The file comprises a number of sections, each starting
  with a single number, which is the number of lines to read in the
  next section. These lines have two numbers each, the $x$ and $y$
  coordinates of a vertex of the map outline.
\end{description}

% \section{Reading in and displaying data}
% \label{manual:sec:datafile-utils}

% All code is to be found in the \texttt{trunk/src} directory. The R
% program should be started from this directory in the following examples.

% To read in data use the functions in \texttt{datafile-utils.R}. Here
% is code to read in the data from a directory containing the
% \texttt{SYS.SYS} and \texttt{ALU.MAP} files, as detailed above:
% \begin{verbatim}
% source("datafile-utils.R")
% sys <- read.sys("/data/path/gm257-1-P8-C57BL6/")
% map <- read.map("/data/path/gm257-1-P8-C57BL6/")
% plot.sys.map(sys, map)
% \end{verbatim}

\begin{thebibliography}{99}
\bibitem{SterrattEtal2012} Sterratt DC, Lyngholm D, Willshaw DJ,
  Thompson ID (2013) Standard anatomical and visual space for the
  mouse retina: computational reconstruction and transformation of
  flattened retinae with the Retistruct package \newblock PLoS
  Comp. Biol. 9
\end{thebibliography}

\end{document}

%%% Local Variables: 
%%% TeX-PDF-mode: t
%%% End: 

% LocalWords:  MAPNUM MINLAT MAXLAT MINLON MAXLON LABLAT LABLON XRED YRED XGRID
% LocalWords:  XGREEN YGREEN XDOUBLE labeled YGRID PERIM GRIDX GRIDY XGRIDCOO
% LocalWords:  YGRIDCOO BOXSIZEX BOXSIZEY TOTALCEL TOTALRED TOTALGRE TOTALDOU
% LocalWords:  MEANPERI MEANAREA SYS SYSTAT ALU src datafile utils myheadings
% LocalWords:  Ubuntu retistruct Datapoints PDF csv Dataset nflip EOD YDOUBLE
% LocalWords:  Matlab cd matlab lX Dss Tss Sss IDT ImageJ png Analyze roi DV xs
% LocalWords:  Metadata metadata datapoints polararea DssMean KDE ys Xcode App
%  LocalWords:  Ctrl gui OpenOffice Calc Karcher DssHullarea flevels Xquartz KR
%  LocalWords:  maxs Lattitude gpa Sterratt Lyngholm Willshaw PLoS GTK retinae
% LocalWords:  alt isn checkboxes Markup markup toolbar min max von
%  LocalWords:  Csefalvay's workspace od ReconstructedOutline getDss
%  LocalWords:  ReconstructedDataset getSss
